\documentclass[12pt, a4paper]{article}
\usepackage[T2A]{fontenc}
\usepackage[utf8]{inputenc}
\usepackage[russian]{babel}
\usepackage{amsmath, amssymb, amsthm}
\usepackage{geometry}
\usepackage{bm}
\geometry{left=2.5cm, right=2cm, top=2cm, bottom=2cm}

\theoremstyle{definition}
\newtheorem*{theory}{Теория}

\begin{document}

\begin{center}
    \Large\bfseries Задача №9: Решение систем нелинейных уравнений методом Ньютона с модификацией Федоренко \\
\end{center}

\vspace{0.5cm}

\section{Введение}
В данной работе рассматривается задача численного поиска корней системы нелинейных уравнений $F(\bm{u}) = 0$. Используется классический многомерный метод Ньютона-Канторовича, дополненный эвристической модификацией Федоренко. Эта модификация включает в себя процедуру дробления шага и введение специальной взвешенной нормы невязки для обеспечения глобальной сходимости метода даже при неудачном начальном приближении.

В качестве тестового примера решается система нелинейных уравнений:
\[
\begin{cases} 
f_1(x,y) = x^5 + y^4 - 2 = 0 \\ 
f_2(x,y) = (x-2)^3 + (y-3)^3 + 16 = 0 
\end{cases}
\]

\section{Теоретическая база: Метод Ньютона-Канторовича}
\begin{theory}
Пусть задана система нелинейных уравнений $F(\bm{u}) = 0$, где $\bm{u} = (x, y)^T$.
Разложим вектор-функцию $F(\bm{u})$ в многомерный ряд Тейлора в окрестности текущего приближения $\bm{u}_n$:
\[
F(\bm{u}) \approx F(\bm{u}_n) + J(\bm{u}_n) (\bm{u} - \bm{u}_n)
\]
где $J(\bm{u}_n) = \frac{\partial F}{\partial \bm{u}}(\bm{u}_n)$ --- матрица Якоби системы, вычисленная в точке $\bm{u}_n$.

Приравнивая линейную часть к нулю, получаем систему линейных алгебраических уравнений (СЛАУ) относительно поправки $\Delta \bm{u}_n = \bm{u}_{n+1} - \bm{u}_n$:
\[
J(\bm{u}_n) \Delta \bm{u}_n = -F(\bm{u}_n)
\]
\end{theory}

Для системы размерности $2 \times 2$ матрица Якоби имеет вид:
\[
J(x, y) = \begin{pmatrix} 
\frac{\partial f_1}{\partial x} & \frac{\partial f_1}{\partial y} \\ 
\frac{\partial f_2}{\partial x} & \frac{\partial f_2}{\partial y} 
\end{pmatrix}
\]
Тогда вектор поправки $\Delta \bm{u}_n = (\Delta x, \Delta y)^T$ можно найти аналитически по формулам Крамера:
\begin{align*}
\Delta x &= \frac{1}{\det J} \left( -f_1 \frac{\partial f_2}{\partial y} + f_2 \frac{\partial f_1}{\partial y} \right) \\
\Delta y &= \frac{1}{\det J} \left( -f_2 \frac{\partial f_1}{\partial x} + f_1 \frac{\partial f_2}{\partial x} \right)
\end{align*}
где определитель $\det J = \frac{\partial f_1}{\partial x} \frac{\partial f_2}{\partial y} - \frac{\partial f_1}{\partial y} \frac{\partial f_2}{\partial x}$.

Стандартный итерационный шаг метода Ньютона выглядит как $\bm{u}_{n+1} = \bm{u}_n + \Delta \bm{u}_n$. Однако этот метод обладает лишь \textit{локальной} сходимостью и может расходиться, если начальное приближение $\bm{u}_0$ выбрано далеко от корня.

\section{Модификация Федоренко (Глобализация сходимости)}

Для расширения области сходимости метода вводится параметр дробления шага (релаксации) $\tau \in (0, 1]$. Новая точка вычисляется как:
\[
\bm{u}_{n+1} = \bm{u}_n + \tau \Delta \bm{u}_n
\]

\subsection{Взвешенная норма невязки}
Обычная евклидова норма невязки $\|F(\bm{u}_n)\|_2$ может вести себя непредсказуемо, особенно если функции в системе имеют существенно разную скорость роста (в нашем случае полиномы 5-й и 3-й степени). Для корректного контроля сходимости вводится \textbf{взвешенная норма}. 

Для каждой функции $f_i(x,y)$ вычисляется локальный коэффициент нормировки (норма градиента):
\[
\alpha_i(x, y) = \|\nabla f_i(x, y)\|_2 = \sqrt{ \left(\frac{\partial f_i}{\partial x}\right)^2 + \left(\frac{\partial f_i}{\partial y}\right)^2 }
\]
Тогда общая взвешенная норма невязки системы определяется как:
\[
\|F(\bm{u}_n)\|_{w} = \sqrt{ \sum_{i=1}^2 \left( \frac{f_i(\bm{u}_n)}{\alpha_i(\bm{u}_n)} \right)^2 }
\]
Геометрический смысл такой нормировки: мы делим значение функции на тангенс угла наклона её поверхности, получая оценку реального геометрического расстояния до поверхности нулевого уровня.

\subsection{Алгоритм выбора шага $\tau$}
На каждой итерации $n$ начальное значение шага полагается равным $\tau = 1$ (полный ньютоновский шаг). Проверяется условие монотонного убывания взвешенной невязки:
\[
\|F(\bm{u}_n + \tau \Delta \bm{u}_n)\|_{w} < \|F(\bm{u}_n)\|_{w}
\]
Если условие не выполняется, шаг дробится вдвое ($\tau \to \tau / 2$) до тех пор, пока не будет достигнуто убывание невязки. Это гарантирует, что метод не "выскочит" из области притяжения корня и будет стабильно сходиться к решению.

\section{Реализация и результаты}
Алгоритм реализован на C++. Вычисление частных производных для формирования якобиана закодировано аналитически:
\begin{align*}
\frac{\partial f_1}{\partial x} &= 5x^4, \quad \frac{\partial f_1}{\partial y} = 4y^3 \\
\frac{\partial f_2}{\partial x} &= 3(x-2)^2, \quad \frac{\partial f_2}{\partial y} = 3(y-3)^2
\end{align*}

В ходе численного эксперимента при старте из начального приближения $x_0 = 1.0, y_0 = 1.0$ программа успешно находит корень за \textbf{13 итераций}. Получены следующие численные значения:
\begin{itemize}
    \item Найденное решение: $x^* \approx 1.14057, \quad y^* \approx 0.513933$
    \item Значения функций в точке оптимума: $f_1 \approx -2.21 \times 10^{-10}, \quad f_2 \approx 5.32 \times 10^{-12}$
    \item Финальная евклидова ошибка $\|F\|_2$: $\approx 2.21 \times 10^{-10}$
\end{itemize}

Критерий останова по взвешенной невязке ($\|F\|_w < 10^{-10}$) уверенно достигается. Модификация Федоренко эффективно стабилизирует метод, обеспечивая надежную сходимость и предотвращая «вылеты» в область бесконечности на начальных итерациях, что критически важно при работе с полиномами высоких (четвертой и пятой) степеней.

\end{document}